\section{The People of Words}

An important aspect to Tradition --- which is difficult for those committed to egalitarianism to accept --- is that not everyone is constitutionally able to grasp the principles of metaphysics. That is why there cannot be a ``traditionalist school'', or system of thought, open to everyone willy-nilly. That is why there can be no question of ``influence''. There are those who know and those who believe; knowledge is infallible by definition. Beliefs are open to debate. \textbf{Rene Guenon} explains in \emph{Spiritual Authority and Temporal Power}.

\begin{quotex}
Knowledge par excellence --- the only kind that truly deserves that name in its fullest sense --- is knowledge of principles, independent of all contingent applications and this belongs exclusively to those who possess spiritual authority because there is nothing in it deriving from the temporal order.
\end{quotex}


The most difficult notion to get across is that words are at best the means, not the goal; rather the goal is metaphysical realization. \textbf{Ibn Arabi}, in his meditation on Seth mentions the \emph{mutakallim}, i.e., the people of words. These are the theologians and philosophers who are unable to practice reflection and meditation. Ibn Arabi describes such people in the meditation on Jethro:

\begin{quotex}
As for the philosophers and masters of thought among the ancients and the mutakallimun in their discourse on the self and its whatness, none of them stumbled onto its reality, and logical speculation can never provide it. He who seeks knowledge of it by means of logical speculation, swells himself up and boasts without vigour or substance. ``They are those whose efforts in the life of this world are misguided while they suppose that they are doing good.'' (Quran 18:104) He who seeks the matter by other than its path will not achieve its realization.
\end{quotex}


It is hubris to presume that God can be known just by reading some books or following the ``perennial philosophy'', as though God can be encompassed in some philosophical system. Any profane philosophy, even so-called perennial philosophy, is about the temporal order, the world of the senses. As such, it can never achieve true knowledge.

Nor does Tradition have anything to do with the study of ancient religions. Tradition means what is handed down; it does not mean a return to some imagined past. The Infinite cannot repeat itself, so there is no going back.

Esoteric knowledge is not some hidden ``secret'' that you share with your friends, that supposedly provides a key to understand everything. It is an ``open'' secret in the sense that it is in public view. However, some will be able to see it, but most will not. As Guenon writes:

\begin{quotationx}
truths of a certain order by their very nature resist all `popularization': however clearly that they are set out and only those who are qualified to understand them will understand them, and for all others they will be as if they did not exist.

In the end, the real secret, the only secret than can never be betrayed in any way, resides uniquely in the inexpressible, which is by the same token incommunicable, every truth of a transcendent order necessarily partaking of the inexpressible
\end{quotationx}

Most decidedly, it is not about a ``school'' that enters into debate with other schools. The object of metaphysics is being, and that means the knowledge of essences. Knowledge is infallible so as soon as this truth is realized, there is no need to get involved in debates.

\paragraph{Knowledge}

\begin{quotex} 
The \emph{mutakallim} sets himself apart from the realised ones among the People of God, who are endowed with unveiling and real existence.
\flright{\textsc{Ibn Arabi}}
\end{quotex}

Knowledge begins with the senses, and then the imagination; symbols are in the imagination. However, it is necessary to abstract from the images to reach the pure intellect. Ultimately, even thoughts must be left behind.

It is said that when you know yourself, you know God. That is not because of human effort. Rather, God is the mirror in which you see yourself in your perfection because you can never see God, just as you don't see the mirror when you look at your reflection.

Our goal is therefore to await the unveiling of God and real existence, not the accumulation of words. If that is our disposition, then the prayer will be answered.

\paragraph{Note}

This is as it should be, since not everyone has the aptitude nor the interest in the subject. Most people are more interested in practical matters, i.e., activities in the temporal order. So the point of this post is to address the abuse of metaphysical ideas.

\flrightit{Posted on 2021-09-05 by Cologero}

